% GENERATED FILE. Edit canonical lessons, then run npm run generate:books.
% canonical-source-hash: fnv1a64:7fa77f14a6210b6d
% canonical-lessons: PA-C239-tava, PA-C239-sekna, PA-C239-chimta, PA-C239-dhakkna, PA-C239-ghiu

\chapter{Bread on the Griddle}
\label{ch:pa-a2-bread-on-the-griddle}
\hlchaptermodality{239}

\begin{chapteropening}
\textbf{By the end of this chapter:} \emph{I can talk about cooking flatbread.}

Complete the last lesson of chapter 239: I can talk about cooking flatbread.
\end{chapteropening}

\section[\texorpdfstring{\pa{ਤਵਾ} (tavā)}{ਤਵਾ (tavā)}]{\pa{ਤਵਾ} (tavā) — a griddle}

\label{lesson:PA-C239-tava}



\begin{quote}
Before the new one: say the Punjabi for a ball of dough, then the Punjabi for to roll out.
\end{quote}

\subsection*{You'll want to know: \pa{ਤਵਾ}}
\textbf{\pa{ਤਵਾ}} — \emph{tavā} — \textquotedblleft{}a griddle\textquotedblright{}.

\begin{quote}
\pa{ਤਵਾ} \pa{ਚੁੱਲ੍ਹੇ} \pa{ਉੱਤੇ} \pa{ਹੈ।}
\end{quote}

\begin{quote}
\emph{tavā chullhe utte hai.}
\end{quote}

\begin{quote}
The griddle is on the stove.
\end{quote}

A masculine noun.

\begin{cousinweb}
\textbf{\pa{ਤਵਾ}} (\emph{tavā}) is inherited from a reconstructed \emph{*tapaka}, \textquotedblleft{}frying pan\textquotedblright{}; CDIAL lists Punjabi \emph{tavā}, \textquotedblleft{}circular griddle\textquotedblright{} (CDIAL 5670).
\end{cousinweb}

\begin{grammarlens}[title={What you've built}]
One more word for this chapter.
\end{grammarlens}

\subsection*{Your turn}
\begin{itemize}
\raggedright
  \item \emph{Say it:} \emph{tavā}
  \item \emph{Say it:} \emph{tavā}, once more
  \item \emph{Say it:} \emph{tavā chullhe utte hai.} — The griddle is on the stove.
  \item \emph{Recall:} say the Punjabi for a ball of dough, then the Punjabi for to roll out, then say \emph{tavā} again
\end{itemize}

\begin{tcolorbox}[breakable,colback=teal!4,colframe=teal!35!black,title={Before you move on}]
What does \pa{ਤਵਾ} mean? (\textquotedblleft{}A griddle\textquotedblright{}.) Say it once more.
\end{tcolorbox}
\section[\texorpdfstring{\pa{ਸੇਕਣਾ} (sekṇā)}{ਸੇਕਣਾ (sekṇā)}]{\pa{ਸੇਕਣਾ} (sekṇā) — to bake on a griddle; to warm}

\label{lesson:PA-C239-sekna}



\begin{quote}
Before the new one: say the Punjabi for to roll out, then the Punjabi for a griddle.
\end{quote}

\subsection*{You'll want to know: \pa{ਸੇਕਣਾ}}
\textbf{\pa{ਸੇਕਣਾ}} — \emph{sekṇā} — \textquotedblleft{}to bake on a griddle; to warm\textquotedblright{}.

\begin{quote}
\pa{ਮੈਂ} \pa{ਤਵੇ} \pa{ਉੱਤੇ} \pa{ਰੋਟੀ} \pa{ਸੇਕਦੀ} \pa{ਹਾਂ।}
\end{quote}

\begin{quote}
\emph{maiṁ tave utte roṭī sekdī hāṁ.}
\end{quote}

\begin{quote}
I bake the roti on the griddle.
\end{quote}

Past (I …): \pa{ਮੈਂ} \pa{ਰੋਟੀ} \pa{ਸੇਕੀ}, \emph{maiṁ roṭī sekī}.

\begin{cousinweb}
\textbf{\pa{ਸੇਕਣਾ}} (\emph{sekṇā}) is inherited from a reconstructed \emph{*sēkk}, \textquotedblleft{}heat, foment, roast\textquotedblright{}; CDIAL lists Punjabi \emph{sekṇā}, \textquotedblleft{}to warm, toast, foment\textquotedblright{} (CDIAL 13581).
\end{cousinweb}

\begin{grammarlens}[title={What you've built}]
One more word for this chapter.
\end{grammarlens}

\subsection*{Your turn}
\begin{itemize}
\raggedright
  \item \emph{Say it:} \emph{sekṇā}
  \item \emph{Say it:} \emph{sekṇā}, once more
  \item \emph{Say it:} \emph{maiṁ tave utte roṭī sekdī hāṁ.} — I bake the roti on the griddle.
  \item \emph{Recall:} say the Punjabi for to roll out, then the Punjabi for a griddle, then say \emph{sekṇā} again
  \item \emph{Write it:} copy \textbf{\pa{ਸੇਕਣਾ}} once, checking the whole word
\end{itemize}

\begin{tcolorbox}[breakable,colback=teal!4,colframe=teal!35!black,title={Before you move on}]
What does \pa{ਸੇਕਣਾ} mean? (\textquotedblleft{}To bake on a griddle; to warm\textquotedblright{}.) Say it once more.
\end{tcolorbox}
\section[\texorpdfstring{\pa{ਚਿਮਟਾ} (chimṭā)}{ਚਿਮਟਾ (chimṭā)}]{\pa{ਚਿਮਟਾ} (chimṭā) — tongs}

\label{lesson:PA-C239-chimta}



\begin{quote}
Before the new one: say the Punjabi for a griddle, then the Punjabi for to bake on a griddle.
\end{quote}

\subsection*{You'll want to know: \pa{ਚਿਮਟਾ}}
\textbf{\pa{ਚਿਮਟਾ}} — \emph{chimṭā} — \textquotedblleft{}tongs\textquotedblright{}.

\begin{quote}
\pa{ਮੈਂ} \pa{ਚਿਮਟੇ} \pa{ਨਾਲ} \pa{ਰੋਟੀ} \pa{ਚੁੱਕਦੀ} \pa{ਹਾਂ।}
\end{quote}

\begin{quote}
\emph{maiṁ chimṭe nāl roṭī chukkdī hāṁ.}
\end{quote}

\begin{quote}
I lift the roti with tongs.
\end{quote}

A masculine noun.

\begin{cousinweb}
\textbf{\pa{ਚਿਮਟਾ}} (\emph{chimṭā}) is inherited from a reconstructed \emph{*cimb}, \textquotedblleft{}pinch\textquotedblright{}; CDIAL lists Punjabi \emph{cimṭā}, \textquotedblleft{}tongs\textquotedblright{} (CDIAL 4822).
\end{cousinweb}

\begin{grammarlens}[title={What you've built}]
One more word for this chapter.
\end{grammarlens}

\subsection*{Your turn}
\begin{itemize}
\raggedright
  \item \emph{Say it:} \emph{chimṭā}
  \item \emph{Say it:} \emph{chimṭā}, once more
  \item \emph{Say it:} \emph{maiṁ chimṭe nāl roṭī chukkdī hāṁ.} — I lift the roti with tongs.
  \item \emph{Recall:} say the Punjabi for a griddle, then the Punjabi for to bake on a griddle, then say \emph{chimṭā} again
  \item \emph{Write it:} copy \textbf{\pa{ਚਿਮਟਾ}} once, checking the whole word
\end{itemize}

\begin{tcolorbox}[breakable,colback=teal!4,colframe=teal!35!black,title={Before you move on}]
What does \pa{ਚਿਮਟਾ} mean? (\textquotedblleft{}Tongs\textquotedblright{}.) Say it once more.
\end{tcolorbox}
\section[\texorpdfstring{\pa{ਢੱਕਣਾ} (ḍhakkṇā)}{ਢੱਕਣਾ (ḍhakkṇā)}]{\pa{ਢੱਕਣਾ} (ḍhakkṇā) — to cover}

\label{lesson:PA-C239-dhakkna}



\begin{quote}
Before the new one: say the Punjabi for to bake on a griddle, then the Punjabi for tongs.
\end{quote}

\subsection*{You'll want to know: \pa{ਢੱਕਣਾ}}
\textbf{\pa{ਢੱਕਣਾ}} — \emph{ḍhakkṇā} — \textquotedblleft{}to cover\textquotedblright{}.

\begin{quote}
\pa{ਮੈਂ} \pa{ਰੋਟੀਆਂ} \pa{ਕੱਪੜੇ} \pa{ਨਾਲ} \pa{ਢੱਕਦੀ} \pa{ਹਾਂ।}
\end{quote}

\begin{quote}
\emph{maiṁ roṭīāṁ kappṛe nāl ḍhakkdī hāṁ.}
\end{quote}

\begin{quote}
I cover the rotis with a cloth.
\end{quote}

Past (I …): \pa{ਮੈਂ} \pa{ਪਤੀਲਾ} \pa{ਢੱਕਿਆ}, \emph{maiṁ patīlā ḍhakkiā}.

\begin{cousinweb}
\textbf{\pa{ਢੱਕਣਾ}} (\emph{ḍhakkṇā}) is inherited from a reconstructed \emph{*ḍhakk}, \textquotedblleft{}cover\textquotedblright{}; CDIAL lists Punjabi \emph{ḍhakkṇā}, \textquotedblleft{}to cover\textquotedblright{} (CDIAL 5574).
\end{cousinweb}

\begin{grammarlens}[title={What you've built}]
One more word for this chapter.
\end{grammarlens}

\subsection*{Your turn}
\begin{itemize}
\raggedright
  \item \emph{Say it:} \emph{ḍhakkṇā}
  \item \emph{Say it:} \emph{ḍhakkṇā}, once more
  \item \emph{Say it:} \emph{maiṁ roṭīāṁ kappṛe nāl ḍhakkdī hāṁ.} — I cover the rotis with a cloth.
  \item \emph{Recall:} say the Punjabi for to bake on a griddle, then the Punjabi for tongs, then say \emph{ḍhakkṇā} again
  \item \emph{Write it:} copy \textbf{\pa{ਢੱਕਣਾ}} once, checking the whole word
\end{itemize}

\begin{tcolorbox}[breakable,colback=teal!4,colframe=teal!35!black,title={Before you move on}]
What does \pa{ਢੱਕਣਾ} mean? (\textquotedblleft{}To cover\textquotedblright{}.) Say it once more.
\end{tcolorbox}
\section[\texorpdfstring{\pa{ਘਿਉ} (ghiu)}{ਘਿਉ (ghiu)}]{\pa{ਘਿਉ} (ghiu) — ghee; clarified butter}

\label{lesson:PA-C239-ghiu}



\begin{quote}
Before the new one: say the Punjabi for tongs, then the Punjabi for to cover.
\end{quote}

\subsection*{You'll want to know: \pa{ਘਿਉ}}
\textbf{\pa{ਘਿਉ}} — \emph{ghiu} — \textquotedblleft{}ghee; clarified butter\textquotedblright{}.

\begin{quote}
\pa{ਦਾਦੀ} \pa{ਰੋਟੀ} \pa{ਉੱਤੇ} \pa{ਘਿਉ} \pa{ਮਲਦੀ} \pa{ਹੈ।}
\end{quote}

\begin{quote}
\emph{dādī roṭī utte ghiu maldī hai.}
\end{quote}

\begin{quote}
Grandmother rubs ghee on the roti.
\end{quote}

A masculine noun.

\begin{cousinweb}
\textbf{\pa{ਘਿਉ}} (\emph{ghiu}) is inherited from Sanskrit \emph{ghṛta}, \textquotedblleft{}clarified butter, ghee\textquotedblright{}; CDIAL lists Punjabi \emph{gheu}, \textquotedblleft{}ghee\textquotedblright{} (CDIAL 4501).
\end{cousinweb}

\begin{grammarlens}[title={What you've built}]
One more word for this chapter.
\end{grammarlens}

\subsection*{Your turn}
\begin{itemize}
\raggedright
  \item \emph{Say it:} \emph{ghiu}
  \item \emph{Say it:} \emph{ghiu}, once more
  \item \emph{Say it:} \emph{dādī roṭī utte ghiu maldī hai.} — Grandmother rubs ghee on the roti.
  \item \emph{Recall:} say the Punjabi for tongs, then the Punjabi for to cover, then say \emph{ghiu} again
  \item \emph{Write it:} copy \textbf{\pa{ਘਿਉ}} once, checking the whole word
\end{itemize}

\begin{tcolorbox}[breakable,colback=teal!4,colframe=teal!35!black,title={Before you move on}]
What does \pa{ਘਿਉ} mean? (\textquotedblleft{}Ghee; clarified butter\textquotedblright{}.) Say it once more.
\end{tcolorbox}
